\section{Introduction}
\label{sec:introduction}
Code-generating assistants introduce an interaction decision alongside the
technical task of producing useful code: assistance can appear automatically or
be requested explicitly. Understanding this distinction requires examining how
programmers use suggestions, rather than assuming that access to a model alone
determines productivity. Barke et al.\ describe interaction patterns in which
programmers use Copilot for both acceleration and exploration~\cite{barke2023grounded}.
Ziegler et al.\ connect perceived productivity with suggestion acceptance,
illustrating that usage measures and productivity judgments need careful
interpretation~\cite{ziegler2022productivity}.

Our starting point is the published study of Copilot interaction modes by
Guglielmi et al.~\cite{guglielmi2026copilot}. We do not claim that the relationship
between experience and suggestion mode is an unexplored topic. Instead, the
question for this replication is whether a transparent participant-level
analysis of the released productivity tables supports an interpretable comparison
of mode-related changes across experience groups. A difference between two
within-group significance decisions would not itself establish a difference
between the groups. We therefore estimate that contrast directly.

We use automatic suggestions as the comparison condition and manual Copilot
invocation as the baseline. This is a comparison of interaction modes, not AI
assistance versus unaided programming. Our exploratory approach combines source
validation, paired productivity changes, and uncertainty estimates. It does not
reproduce every outcome or model in the source study.

\subsection{Contributions and Research Questions}
\label{sec:questions}
The contributions of this course project are:
\begin{itemize}
  \item a reproducible validation pipeline for the published productivity tables,
  including checksum verification and independent metric recalculation;
  \item explicit participant-level estimands for within-group changes and the
  expert-minus-novice contrast, retaining the paired design;
  \item joint reporting of productivity, completion, and time, with documented
  limits on what these measurements can establish.
\end{itemize}
We ask three research questions:
\begin{enumerate}
  \item[RQ1:] What is the mean paired productivity change from manual to automatic
  suggestions among novices?
  \item[RQ2:] What is the corresponding change among experts?
  \item[RQ3:] How does the mean paired change differ between experts and novices?
\end{enumerate}
RQ3 is primary; RQ1 and RQ2 provide context. Section~\ref{sec:methodology}
defines the dataset and estimands. Section~\ref{sec:results} identifies the
remaining results integration, and Section~\ref{sec:limitations} discusses validity.
